\noindent W-bosons predominantly decay into quark-antiquark pairs, but unfortunately this decay mode is not distinguishable from the general QCD background, and the measurement of jet directions and energies is not exact. Hence, we will instead use the decay of the W-boson into lepton and lepton-neutrino pairs. In these $W \rightarrow e \nu_e$ and $W \rightarrow \mu \nu_{\mu}$ decays, the leptons can be easily detected in ATLAS, however this is not true for the neutrinos whose presence and properties must be observed indirectly by the imbalance of the transverse momentum.\\

\noindent When the W-boson decay occurs, in its rest frame, the angular distribution of the leptons with reference to the axis of the incoming partons is given by
$$\frac{\mathrm{d}\sigma}{\mathrm{d}\cos\theta} = 1 + \cos^2\theta$$
By doing a transformation of variables, we can obtain a distribution of the transverse momentum:
$$\frac{\mathrm{d}\sigma}{\mathrm{d}p_{\text{T}}} = \frac{\mathrm{d}\sigma}{\mathrm{d}\cos\theta} \left|\frac{\mathrm{d}p_{\text{T}}}{\mathrm{d}\cos\theta}\right|^{-1} = \frac{\mathrm{d}\sigma}{\mathrm{d}\cos\theta} \frac{2p_{\text{T}}}{m_W} \frac{1}{\sqrt{\frac{1}{4}m_W^2 - p_{\text{T}}^2}}$$
In the lab frame, there will simply be a Lorentz boost along the beam axis and hence the transverse components will remain unchanged. So, what we would observe from the above expression is the existence of a pole at $p_{\text{T}} = m_W/2$, which is called a Jacobi peak. Measuring the location of the Jacobi peak thus lets us calculate the mass of the W-boson.\\

\noindent Ideally, the peak would be at a single point, however in practice it is smeared due to the finite resolution of the detector, the W-decay width, and the transverse momentum $p_{\text{T}}^W$ of the W-boson. This can be dealt with to some extent by measuring instead the location of the half-maximum of the peak. Since the smearing effects affect the distribution symmetrically, the position ofthe half-maximum will remain unchanged.\\